\section{Conclusion}
We introduce FunQA, a challenging VideoQA dataset of counter-intuitive and fun videos designed to explore the depth of video reasoning. We established a benchmark of four rigorous tasks and performed a comprehensive evaluation of frontier baselines, which were revealed to demonstrate a significant performance gap compared to manual annotation. Most baselines struggle especially in tasks requiring higher-level reasoning, such as counter-intuitive timestamp localization. We proposed an agent for VLM: FunMentor. It effectively enhances the model's understanding of humor, demonstrating the powerful potential of agent-based finetuning. FunQA dataset and novel benchmark tasks provide valuable insight and future research direction. We expect FunQA to underpin the next generation of VideoQA research that pushes video reasoning beyond superficial descriptions toward deeper understanding. 